% 2009 年新课标卷理科第 17 题：飞机在 A、B 两点测量两山顶 M、N 的俯角。
% 来源 PDF：普通高考/2009/2009新课标理(海南,宁夏).pdf，
% SHA256 f42b575218180e261a89cf8733109d4960b61560a97b5f266731a51c422a4045
% 原页：第 1 页（A3 横向）右栏下部，印刷页 541。
% 重绘依据：A、B 在同一水平线上，M、N 为下方两山顶；A 处对 M、N 的俯角记为 α、γ，
% B 处对 M 的俯角记为 β、对 N 的俯角记为 δ，AB=d。原图为写意山形，此处按同一几何
% 关系重画并标注一种可行方案所需的数据；方案不唯一。
\begin{tikzpicture}[scale=0.88, semithick, >=Stealth,
  every node/.style={inner sep=1.5pt, font=\small}]
  % 水平基准线
  \draw[dashed] (-0.9,0) -- (10.1,0);
  % 山形
  \draw (0.5,-3.3) -- (1.5,-3.4) -- (2.6,-2.8) -- (3.7,-3.35) -- (4.4,-3.45);
  \draw (6.4,-3.4) -- (7.2,-3.3) -- (8.3,-2.4) -- (9.4,-3.35) -- (9.9,-3.4);
  % 视线
  \draw (0,0) -- (2.6,-2.8);
  \draw (0,0) -- (8.3,-2.4);
  \draw (6,0) -- (2.6,-2.8);
  \draw (6,0) -- (8.3,-2.4);
  % A 处俯角
  \draw (0.85,0) arc (0:-47.1:0.85);
  \node at (1.05,-0.52) {$\alpha$};
  \draw (1.75,0) arc (0:-16.1:1.75);
  \node at (2.05,-0.30) {$\gamma$};
  % B 处俯角
  \draw (5.15,0) arc (180:219.5:0.85);
  \node at (4.75,-0.48) {$\beta$};
  \draw (7.0,0) arc (0:-46.2:1.0);
  \node at (7.10,-0.65) {$\delta$};
  % 点与标注
  \fill (0,0) circle (2.2pt) node[above] {$A$};
  \fill (6,0) circle (2.2pt) node[above] {$B$};
  \fill (2.6,-2.8) circle (2.0pt) node[below] {$M$};
  \fill (8.3,-2.4) circle (2.0pt) node[below] {$N$};
  \draw[<->] (0,0.7) -- (6,0.7);
  \node[above] at (3,0.7) {$d$};
\end{tikzpicture}
